%!TEX root = ../main.tex
% APÉNDICE A — Encuesta a personas cuidadoras

\chapter{Encuesta a personas cuidadoras}\label{apx:encuesta}

\section{Objetivo}

La encuesta buscó conocer cómo organizan hoy el cuidado las personas que
acompañan a alguien con discapacidad, qué les cuesta más y si una
herramienta como la propuesta les resultaría útil. Sus resultados
orientaron las decisiones de diseño descriptas en el
capítulo~\ref{cap:resultados}.

\section{Metodología}

Se elaboró un formulario digital de doce preguntas, compuesto principalmente por preguntas cerradas, dos de respuesta múltiple y una abierta. Antes de su utilización, se realizó una prueba con un grupo de personas cuidadoras, lo que permitió revisar las preguntas y ajustar el formulario. Los resultados se utilizaron como referencia para identificar tendencias y orientar las decisiones tomadas durante el diseño del sistema.

Por indicación de la cátedra, a partir de las respuestas de esa prueba se generaron 80 casos que mantienen su distribución, y sobre ellos se calcularon los porcentajes que se presentan en el capítulo~\ref{cap:resultados}.

\section{Cuestionario}

\begin{enumerate}
    \item ¿Cuida actualmente a una persona con discapacidad? (Sí / No)
    \item ¿Cuál es su relación con la persona que cuida? (Madre o padre / Hijo o hija / Otro familiar / Cuidador o cuidadora profesional)
    \item ¿Hace cuánto tiempo que cuida? (Menos de 1 año / Entre 1 y 3 años / Entre 4 y 7 años / Más de 7 años)
    \item ¿Cuántas horas por día dedica al cuidado? (Menos de 4 / Entre 4 y 8 / Entre 8 y 12 / Más de 12)
    \item ¿Qué tipo de discapacidad tiene la persona? (Motora / Intelectual / Visual / Autismo (TEA) / Múltiple / Otra)
    \item ¿Qué tan difícil le resulta organizar la información relacionada con el cuidado, como turnos, medicación, estudios o certificados? (Muy difícil / Difícil / Ni fácil ni difícil / Fácil)
    \item ¿Cuáles son los problemas más frecuentes que enfrenta? (respuesta múltiple: recordar turnos médicos / controlar la medicación / guardar documentación importante / encontrar información confiable / coordinar tareas con otros cuidadores / sentirse acompañado o recibir orientación)
    \item ¿Qué herramientas usa hoy para organizar el cuidado? (respuesta múltiple: agenda en papel / calendario del celular / WhatsApp / grupo de WhatsApp con la familia / Excel o documentos digitales / ninguna)
    \item En una escala de 1 a 5, ¿cuánto estrés le genera organizar el cuidado?
    \item ¿Una herramienta que reúna toda esta información le sería útil? (Sí, mucho / Sí, bastante / Tal vez / Poco / Nada)
    \item ¿Usaría una aplicación web o móvil para gestionar el cuidado? (Sí / No estoy seguro o segura / No)
    \item ¿Qué es lo que más le ayudaría hoy en la tarea de cuidar? (abierta)
\end{enumerate}

\section{Resumen de la encuesta}

\begin{table}[H]
\centering
\caption{Perfil resumido de los casos}\label{tab:apx-muestra}
\small
\begin{tabularx}{\textwidth}{>{\raggedright\arraybackslash}X r}
\toprule
\tb{Característica} & \tb{Casos} \\
\midrule
Cuidadores profesionales & 50,0~\% \\
Familiares (hijos, madres o padres y otros familiares) & 50,0~\% \\
Cuidan desde hace entre 1 y 3 años & 57,5~\% \\
Dedican entre 4 y 8 horas por día & 60,0~\% \\
Les resulta difícil o muy difícil organizar la información & 57,5~\% \\
Consideran útil una herramienta como la propuesta & 80,0~\% \\
\bottomrule
\end{tabularx}
\end{table}

La pregunta abierta fue respondida por 48 de los 80 casos. Las respuestas se
agruparon por tema: medicación y rutinas, tener la información en un solo
lugar, turnos médicos, acceso al sistema de salud y apoyo y orientación.
